\chapter{Chimie numérique~: Hartree--Fock et DFT}\label{ch:b3:computational-chemistry}

L'ion hydrure d'hélium, \ce{HeH+}, serait la première molécule formée
lors du refroidissement de l'Univers jeune~: un atome d'hélium retenant
un proton. On l'a produit dans un tube à décharge au laboratoire dès
1925, mais on l'a cherché en vain dans l'espace pendant des décennies,
jusqu'à la détection de sa raie de rotation dans une nébuleuse
planétaire en 2019. Sa longueur de liaison, ses fréquences de vibration
et les longueurs d'onde de ses raies étaient alors connues depuis
longtemps --- par le calcul. Un programme qui résout, de façon approchée,
l'\omterm{thm:b3:quantum-model-systems:schrodinger}{équation de Schrödinger} des électrons peut prédire la forme, l'énergie
et le spectre d'une molécule que personne n'a jamais mise en flacon. Ce
chapitre explique comment fonctionnent de tels programmes, de la
séparation des noyaux et des électrons jusqu'à la théorie de
Hartree--Fock et à la \omterm{def:b3:computational-chemistry:dft}{théorie de la fonctionnelle de la densité}, et dit
dans quelle mesure on peut se fier à leurs résultats.

\begin{recall}
Le \Cref{ch:b3:quantum-model-systems} a donné le hamiltonien, l'\omterm{thm:b3:quantum-model-systems:schrodinger}{équation
de Schrödinger} et l'atome d'hydrogène~; le \cref{ch:b3:many-electron-atoms},
le \omterm{def:b3:many-electron-atoms:slater}{déterminant de Slater} et l'\omterm{def:b3:many-electron-atoms:exchange}{intégrale d'échange}. Le volume de deuxième
année a construit les orbitales moléculaires comme combinaisons
linéaires d'orbitales atomiques (CLOA), avec les intégrales de
recouvrement, coulombiennes et de résonance et un déterminant séculaire,
et a appliqué la méthode de Hückel~; il a aussi utilisé les spectres de
photoélectrons pour voir les énergies des orbitales.
\end{recall}

\begin{omfigure}
\includegraphics[width=0.6\linewidth]{images/book4/ai/b3-03-cluster.jpg}

{\small Un centre de calcul~: des rangées de serveurs sur lesquels
tournent aujourd'hui la plupart des calculs de chimie quantique, souvent
pendant des heures ou des jours par molécule.}
\end{omfigure}

\section{Born--Oppenheimer et la surface d'énergie potentielle}

Le hamiltonien d'une molécule contient les énergies cinétiques de ses
noyaux et de ses électrons et toutes leurs interactions coulombiennes.
Les noyaux sont au moins 1836 fois plus lourds que les électrons~: pour
les mêmes forces, ils se déplacent bien plus lentement, et les électrons
s'adaptent presque instantanément à chaque position des noyaux.

\begin{definition}[Approximation de Born--Oppenheimer]\label{def:b3:computational-chemistry:born-oppenheimer}
Dans l'\emph{approximation de Born--Oppenheimer}\index{approximation de Born--Oppenheimer},
on résout l'\omterm{thm:b3:quantum-model-systems:schrodinger}{équation de Schrödinger} électronique pour des noyaux
maintenus fixes aux positions $\mathbf R$,
\[
\hat H_{\mathrm{el}}\,\psi_{\mathrm{el}}(\mathbf r;\mathbf R) =
E_{\mathrm{el}}(\mathbf R)\,\psi_{\mathrm{el}}(\mathbf r;\mathbf R),
\]
puis les noyaux se déplacent dans le potentiel $E_{\mathrm{el}}(\mathbf R) +
V_{\mathrm{nn}}(\mathbf R)$, somme de l'énergie électronique et de la
répulsion des noyaux.
\end{definition}

\begin{definition}[Surface d'énergie potentielle, point stationnaire, optimisation de géométrie]\label{def:b3:computational-chemistry:pes}
La \emph{surface d'énergie potentielle}\index{surface d'energie potentielle@surface d'énergie potentielle} (SEP) d'une
molécule est la fonction $U(\mathbf R) = E_{\mathrm{el}}(\mathbf R) +
V_{\mathrm{nn}}(\mathbf R)$ de ses coordonnées nucléaires. Un \emph{point
stationnaire}\index{point stationnaire} est un point où toutes les
dérivées premières de $U$ s'annulent~; un minimum est une
\emph{géométrie d'équilibre}\index{geometrie d'equilibre@géométrie d'équilibre}.
En chercher une à partir d'une structure de départ, en suivant les
forces $-\nabla U$, est une \emph{optimisation de géométrie}\index{optimisation de geometrie@optimisation de géométrie}.
\end{definition}

Pour une molécule diatomique, la SEP est une courbe $U(R)$~; pour une
molécule triatomique, une fonction de trois coordonnées~; pour $N$ atomes,
de $3N - 6$. La masse des noyaux n'intervient pas dans
$\hat H_{\mathrm{el}}$~: \ce{H2}, \ce{HD} et \ce{D2} ont la
même SEP, la même longueur de liaison et la même \omterm{def:b3:quantum-model-systems:oscillator}{constante de force}, et
ne diffèrent que par la façon dont les noyaux s'y déplacent (\omterm{def:b3:quantum-model-systems:oscillator}{énergie de
point zéro}, \cref{ch:b3:quantum-model-systems}).

\begin{proposition}[Fréquences tirées de la hessienne]\label{prop:b3:computational-chemistry:hessian}
Près d'un minimum, $U \approx U_0 + \frac12\sum_{ij}H_{ij}\,\delta q_i\,\delta q_j$, avec
la hessienne $H_{ij} = \partial^2U/\partial q_i\partial q_j$. En
coordonnées pondérées par les masses $\delta q_i\sqrt{m_i}$, les \omterm{def:b3:quantum-model-systems:operator}{valeurs
propres} $\lambda_k$ de la hessienne pondérée donnent les pulsations
harmoniques $\omega_k =
\sqrt{\lambda_k}$ des modes normaux. En un minimum, tous les $\lambda_k$
sont positifs (six, ou cinq pour une molécule linéaire, sont nuls~:
translations et rotations)~; en un point selle d'ordre un, un seul est
négatif.
\end{proposition}

\begin{proof}[Démonstration partielle]
Pour une molécule diatomique, $U \approx U_0 + \frac12k(R - R_e)^2$ avec
$k = U''(R_e)$, et
l'oscillateur du \cref{ch:b3:quantum-model-systems} donne $\omega =
\sqrt{k/\mu}$. En général, les équations classiques $m_i\ddot q_i =
-\sum_jH_{ij}q_j$ deviennent, en coordonnées pondérées par les masses
$x_i = \sqrt{m_i}q_i$,
$\ddot{\mathbf x} = -\tilde H\mathbf x$ avec la matrice symétrique $\tilde H_{ij} =
H_{ij}/\sqrt{m_im_j}$~; ses vecteurs propres oscillent indépendamment à
$\sqrt{\lambda_k}$. Un $\lambda_k$ négatif donne une fréquence
imaginaire~: dans cette direction, l'énergie diminue, c'est un maximum.
Les modes normaux du \cref{ch:b3:group-theory-applied} sont ces vecteurs
propres.
\end{proof}

\begin{omfigure}
\begin{tikzpicture}[font=\small, scale=0.95]
  \fill[omDef!8] (0,0) ellipse (1.9 and 1.15);
  \foreach \s in {1.0,0.75,0.5,0.25}{\draw[omDef!70] (0,0) ellipse ({1.9*\s} and {1.15*\s});}
  \fill[omThm!8] (6,0.4) ellipse (1.6 and 1.0);
  \foreach \s in {1.0,0.75,0.5,0.25}{\draw[omThm!70] (6,0.4) ellipse ({1.6*\s} and {1.0*\s});}
  \draw[gray] (2.1,-0.15) .. controls (2.9,0.55) and (3.4,0.55) .. (4.3,0.35);
  \draw[gray] (2.2,0.9) .. controls (3.0,0.35) and (3.5,0.35) .. (4.1,1.2);
  \node[omsaddle] at (3.15,0.42) {};
  \draw[ompath] (0,0) .. controls (1.8,0.6) and (2.4,0.5) .. (3.15,0.42) .. controls (4.0,0.35) and (4.8,0.4) .. (6,0.4);
  \node[anchor=north] at (0,-1.25) {minimum A (réactifs)};
  \node[anchor=north] at (6,-0.7) {minimum B (produits)};
  \node[anchor=south] at (3.15,1.25) {point selle};
  \draw[->, gray] (3.15,1.25) -- (3.15,0.6);
  \fill (0,0) circle (1.2pt) (6,0.4) circle (1.2pt);
  \draw[->] (-2.3,-2.2) -- (8.2,-2.2) node[right] {$q_1$};
  \draw[->] (-2.3,-2.2) -- (-2.3,1.8) node[above] {$q_2$};
\end{tikzpicture}

{\small Une \omterm{def:b3:computational-chemistry:pes}{surface d'énergie potentielle} vue de dessus, en courbes de
niveau sur deux coordonnées nucléaires (schéma). Deux minima sont reliés,
par un point selle (croix), par le chemin de plus basse énergie (en
tirets)~: le chemin d'énergie minimale étudié au
\cref{ch:b3:rate-theories}.}
\end{omfigure}

\section{Le principe variationnel}

L'\omterm{thm:b3:quantum-model-systems:schrodinger}{équation de Schrödinger} électronique ne se résout exactement pour
aucune molécule de plus d'un électron. Les approximations se jugent à
l'aune d'un théorème.

\begin{theorem}[Le principe variationnel]\label{thm:b3:computational-chemistry:variational}
Pour toute fonction d'essai normable $\phi$ qui satisfait les conditions
aux limites du problème,
\[
E[\phi] = \frac{\langle\phi|\hat H|\phi\rangle}{\langle\phi|\phi\rangle} \ge E_0,
\]
où $E_0$ est la plus petite \omterm{def:b3:quantum-model-systems:operator}{valeur propre} de $\hat H$~; l'égalité n'a lieu
que si $\phi$ est une \omterm{def:b3:quantum-model-systems:operator}{fonction propre} de l'état fondamental. C'est le
\emph{principe variationnel}\index{principe variationnel}.
\end{theorem}

\begin{proof}
Développons la fonction $\phi = \sum_nc_n\psi_n$ sur les \omterm{def:b3:quantum-model-systems:operator}{fonctions propres} orthonormées de $\hat H$
(\cref{thm:b3:quantum-model-systems:hermitian-real}). Alors $\langle\phi|\hat
H|\phi\rangle = \sum_n|c_n|^2E_n \ge E_0\sum_n|c_n|^2 = E_0\langle\phi|\phi\rangle$,
puisque tout $E_n \ge E_0$~; l'égalité exige $c_n = 0$ dès que $E_n > E_0$.
\end{proof}

Plus l'énergie d'une fonction d'essai est basse, meilleure elle est~;
tout paramètre ajustable se fixe en minimisant $E$.

\begin{proposition}[L'hélium avec une charge écrantée]\label{prop:b3:computational-chemistry:helium-zeta}
Pour l'hélium, la fonction d'essai $\phi = \eu^{-\zeta(r_1 + r_2)}$ (en
unités atomiques~: longueurs en $a_0$, énergies en $E_h$) donne
$E(\zeta) = \zeta^2 - \frac{27}{8}\zeta$,
minimale pour $\zeta = \frac{27}{16}$, avec $E = -(\frac{27}{16})^2 E_h =
\qty{-2.8477}{\hartree}$.
\end{proposition}

\begin{proof}
Pour un électron dans $\eu^{-\zeta r}$, $\langle T\rangle = \zeta^2/2$ et
$\langle 1/r\rangle = \zeta$~; la répulsion de deux électrons dans la même
fonction $1s$ vaut $\frac58\zeta$ (une intégrale classique, admise). Donc $E = 2\cdot\frac12
\zeta^2 - 2\cdot2\zeta + \frac58\zeta = \zeta^2 - \frac{27}{8}\zeta$, et $\dd E/\dd
\zeta = 0$ donne $\zeta = 27/16 = 1.6875$.
\end{proof}

\begin{example}[Quelle est sa qualité~?]\label{ex:b3:computational-chemistry:helium}
% ledger: ie:He, ie4:He2, const:Eh, const:eV
L'énergie totale de l'hélium est l'opposé de la somme de ses deux
énergies d'ionisation,
$-(24.5874 + 54.4178)\,\unit{eV} = \qty{-2.9034}{\hartree}$. La fonction à
un paramètre est trop haute de \qty{0.056}{\hartree} (\qty{1.5}{eV})~:
chaque électron voit un noyau écranté par l'autre jusqu'à
$Z_{\mathrm{eff}} = 1.69$, comme le suggèrent les règles de Slater, mais
les électrons s'évitent aussi d'une manière qu'aucun produit de deux
fonctions ne peut décrire.
\end{example}

\begin{proposition}[Méthode des variations linéaires]\label{prop:b3:computational-chemistry:linear-variation}
Pour une fonction d'essai $\phi = \sum_{i=1}^n c_i\chi_i$ construite sur
$n$ fonctions fixées,
avec $H_{ij} = \langle\chi_i|\hat H|\chi_j\rangle$ et $S_{ij} = \langle\chi_i|\chi_j
\rangle$, les valeurs stationnaires de $E$ sont les racines de
$\det(H - ES) = 0$, et la plus petite racine est un majorant de $E_0$.
\end{proposition}

\begin{proof}
$E\sum_{ij}c_ic_jS_{ij} = \sum_{ij}c_ic_jH_{ij}$ (coefficients réels).
Dérivons par rapport à $c_k$ et posons $\partial E/\partial c_k = 0$~: $\sum_j(H_{kj} -
ES_{kj})c_j = 0$ pour tout $k$. Une solution non nulle exige un
déterminant nul. La plus petite racine est $E[\phi]$ pour son vecteur
propre, donc $\ge E_0$.
\end{proof}

La méthode de Hückel du volume de deuxième année est cette proposition
appliquée à des orbitales $p$ avec des $H_{ij}$ empiriques. Chaque méthode
qui suit est la même idée, avec de meilleures fonctions et de meilleurs
éléments de matrice.

\section{Bases d'orbitales}

\begin{definition}[Base, base minimale, orbitales de type Slater et de type gaussienne]\label{def:b3:computational-chemistry:basis}
Une \emph{base d'orbitales}\index{base d'orbitales} est l'ensemble des fonctions fixées $\chi_i$,
centrées sur les atomes, à partir desquelles on construit les orbitales
moléculaires. Une
\emph{base minimale}\index{base minimale} comporte une fonction par
orbitale atomique occupée (une $1s$ pour H, cinq pour C). Une
\emph{orbitale de type Slater}\index{orbitale de type Slater} (STO) a la forme radiale
$r^{n-1}\eu^{-\zeta r}$~; une \emph{orbitale de type gaussienne}\index{orbitale de type gaussienne}
(GTO) a la forme $\eu^{-\alpha r^2}$ multipliée par un polynôme en $x$, $y$,
$z$. Une \emph{fonction gaussienne contractée}\index{fonction gaussienne contractee@fonction gaussienne contractée}
est une combinaison fixe $\sum_kd_k\,\eu^{-\alpha_kr^2}$ de gaussiennes primitives.
\end{definition}

Les fonctions de Slater ont la bonne forme (un point de rebroussement au
noyau, une queue exponentielle), mais leurs intégrales biélectroniques
sur plusieurs centres sont très coûteuses~; les fonctions gaussiennes
ont la mauvaise forme, mais le produit de deux gaussiennes centrées sur
deux points différents est une gaussienne centrée en un point situé entre
eux, si bien que toute intégrale se ramène à une formule explicite. Les
contractions tirent le meilleur des deux~: dans la base STO-3G, chaque
fonction de Slater est remplacée par une somme fixe de trois gaussiennes
ajustées sur elle.

\begin{example}[La fonction STO-3G de l'hydrogène]\label{ex:b3:computational-chemistry:sto3g}
% ledger: bse:H.sto3g.a1, bse:H.sto3g.a2, bse:H.sto3g.a3, bse:sto3g.c1, bse:sto3g.c2, bse:sto3g.c3
La fonction $1s$ de l'hydrogène en STO-3G est $0.15433\,g(3.42525) + 0.53533\,
g(0.62391) + 0.44463\,g(0.16886)$, où $g(\alpha)$ est la gaussienne normée
$(2\alpha/\pi)^{3/4}\eu^{-\alpha r^2}$ (exposants en $a_0^{-2}$). Elle imite
une fonction de Slater d'exposant $\zeta = 1.24$ --- un atome
d'hydrogène légèrement comprimé, comme il l'est dans les molécules. Son
recouvrement avec la fonction $1s$ exacte de même exposant vaut
$0.9998$~; seuls le point de rebroussement au noyau et la queue lointaine
sont manqués (figure ci-dessous).
\end{example}

\begin{definition}[Fonctions de valence séparée, de polarisation et diffuses]\label{def:b3:computational-chemistry:split-valence}
Une \emph{base à valence séparée}\index{base a valence separee@base à valence séparée} décrit chaque orbitale
de valence par deux fonctions (ou plus) de tailles différentes, de sorte
que l'orbitale puisse grossir ou rétrécir dans une molécule. Une
\emph{fonction de polarisation}\index{fonction de polarisation} a un $l$
plus élevé que les orbitales occupées ($p$ sur H, $d$ sur C) et permet à
la densité de se déplacer hors de l'atome~; une
\emph{fonction diffuse}\index{fonction diffuse} a un petit exposant et
décrit les anions et les interactions faibles loin des noyaux.
\end{definition}

Un nom tel que 6-31G(d) signifie~: six primitives pour chaque fonction de
cœur, la valence séparée en une contraction de trois primitives et une
primitive seule, et des \omterm{def:b3:computational-chemistry:split-valence}{fonctions de polarisation} $d$ sur les atomes
lourds. Les résultats convergent vers la \emph{limite de base complète}
quand la base s'agrandit~; un calcul ne vaut que ce que vaut sa base.

\begin{omfigure}
\begin{tikzpicture}
\begin{axis}[width=0.46\linewidth, height=5.4cm, xmin=0, xmax=5, ymin=0,
  ymax=0.6, axis lines=left, xlabel={$r$ / $a_0$}, ylabel={$\phi(r)$},
  legend cell align=left, legend style={font=\small, draw=none},
  tick label style={font=\small}, label style={font=\small},
  title={\small STO-3G comparée à $\eu^{-r}$ ($\zeta = 1$)}]
\addplot[black, thick] table[x=r, y=slater] {figdata/out/bachelor-3-minimal-basis-hf-sto.dat};
\addlegendentry{Slater}
\addplot[omDef, thick, dashed] table[x=r, y=g3] {figdata/out/bachelor-3-minimal-basis-hf-sto.dat};
\addlegendentry{STO-3G}
\addplot[omThm, thick, dotted] table[x=r, y=g1] {figdata/out/bachelor-3-minimal-basis-hf-sto.dat};
\addlegendentry{une gaussienne}
\end{axis}
\end{tikzpicture}\hfill
\begin{tikzpicture}
\begin{axis}[width=0.46\linewidth, height=5.4cm, xmin=0.5, xmax=8.5,
  ymin=-2.86, ymax=-2.78, axis lines=left, xlabel={cycle SCF},
  ylabel={$E$ / $E_h$}, ytick={-2.85,-2.83,-2.81,-2.79},
  yticklabel style={/pgf/number format/fixed, /pgf/number format/precision=2},
  tick label style={font=\small}, label style={font=\small},
  title={\small SCF de \ce{HeH+}}]
\addplot[omDef, thick, mark=*] table[x=it, y=E] {figdata/out/bachelor-3-minimal-basis-hf-heh-scf.dat};
\addplot[gray, dashed, domain=0.5:8.5] {-2.84184};
\end{axis}
\end{tikzpicture}

{\small À gauche~: la contraction STO-3G (en tirets) reproduit fidèlement
la fonction de Slater $1s$, sauf au noyau~; une gaussienne seule (en
pointillés) n'y parvient pas. À droite~: l'énergie de \ce{HeH+} à
\qty{1.4632}{\bohr} au fil des itérations du \omterm{def:b3:computational-chemistry:hartree-fock}{champ autocohérent} (STO-3G)~;
elle converge par valeurs supérieures vers
\qty{-2.8418}{\hartree} en quelques cycles, chaque itéré obéissant au
\omterm{thm:b3:computational-chemistry:variational}{principe variationnel}.}
\end{omfigure}

\section{La théorie de Hartree--Fock}

La \omterm{def:b3:many-electron-atoms:indistinguishable}{fonction d'onde antisymétrique} la plus simple de $N$ électrons est un
\omterm{def:b3:many-electron-atoms:slater}{déterminant de Slater} unique. La théorie de Hartree--Fock trouve le
meilleur.

\begin{definition}[Méthode de Hartree--Fock, opérateur de Fock, champ autocohérent]\label{def:b3:computational-chemistry:hartree-fock}
La \emph{méthode de Hartree--Fock}\index{methode de Hartree--Fock@méthode de Hartree--Fock} approche
l'état fondamental par le \omterm{def:b3:many-electron-atoms:slater}{déterminant de Slater} unique d'énergie la plus
basse. Ses orbitales sont \omterm{def:b3:quantum-model-systems:operator}{fonctions propres} de l'\emph{opérateur de Fock}\index{operateur de Fock@opérateur de Fock}
\[
\hat f(1) = \hat h(1) + \sum_b\big[2\hat J_b(1) - \hat K_b(1)\big],
\]
dans le cas à couches complètes~: $\hat h$ regroupe l'énergie cinétique
d'un électron et son attraction par les noyaux, $\hat J_b$ la répulsion
par le nuage de charge de l'orbitale $b$, et $\hat K_b$ est l'\omterm{def:b3:quantum-model-systems:operator}{opérateur}
d'échange, dont les \omterm{def:b3:quantum-model-systems:hermitian}{valeurs moyennes} sont les \omterm{def:b3:many-electron-atoms:exchange}{intégrales d'échange} du
\cref{ch:b3:many-electron-atoms}. Comme $\hat f$ dépend des orbitales
qu'il détermine, on résout les équations par itération jusqu'à ce que
les orbitales ne changent plus~: c'est un \emph{champ
autocohérent}\index{champ autocoherent@champ autocohérent} (SCF).
\end{definition}

\begin{theorem}[Équations de Roothaan--Hall]\label{thm:b3:computational-chemistry:roothaan}
En écrivant chaque orbitale dans une base, $\phi_a = \sum_\mu C_{\mu a}\chi_\mu$, les
équations de Hartree--Fock à couches complètes deviennent le problème
matriciel aux \omterm{def:b3:quantum-model-systems:operator}{valeurs propres}
\[
\mathbf F\mathbf C = \mathbf S\mathbf C\boldsymbol\varepsilon, \qquad F_{\mu\nu} =
h_{\mu\nu} + \sum_{\lambda\sigma}P_{\lambda\sigma}\big[(\mu\nu|\sigma\lambda) -
\tfrac12(\mu\lambda|\sigma\nu)\big],
\]
avec la matrice densité $P_{\lambda\sigma} = 2\sum_a^{\mathrm{occ}}C_{\lambda a}C_{\sigma
a}$ et les intégrales biélectroniques $(\mu\nu|\lambda\sigma) = \iint\chi_\mu(1)\chi_\nu(1)
r_{12}^{-1}\chi_\lambda(2)\chi_\sigma(2)$ (unités atomiques).
\end{theorem}

\begin{proof}[Démonstration partielle]
L'énergie du déterminant vaut $E = \sum_{\mu\nu}P_{\mu\nu}h_{\mu\nu} + \frac12
\sum P_{\mu\nu}P_{\lambda\sigma}[(\mu\nu|\sigma\lambda) - \frac12(\mu\lambda|\sigma\nu)]$,
fonction quadratique des coefficients. On la minimise sous les
contraintes que les orbitales restent orthonormées, $\mathbf C^{\mathsf T}\mathbf S\mathbf C = \mathbf 1$,
avec un multiplicateur de Lagrange par contrainte~: les conditions de
stationnarité s'écrivent $\mathbf F\mathbf C = \mathbf S\mathbf C\boldsymbol\varepsilon$,
les multiplicateurs formant une matrice que l'on peut rendre diagonale
par une rotation des orbitales occupées entre elles. Le calcul est admis.
\end{proof}

\begin{method}[La procédure SCF]\label{met:b3:computational-chemistry:scf}
\begin{enumerate}
  \item Calculer une fois pour toutes les intégrales $S_{\mu\nu}$, $h_{\mu\nu}$ et $(\mu\nu|\lambda\sigma)$.
  \item Estimer la matrice densité (par exemple à partir des vecteurs
        propres de $\mathbf h$ seule, l'estimation «~de cœur~»).
  \item Construire $\mathbf F$ à partir de $\mathbf P$~; résoudre $\mathbf F\mathbf C = \mathbf S\mathbf C
        \boldsymbol\varepsilon$ (orthogonaliser avec $\mathbf S^{-1/2}$,
        puis diagonaliser).
  \item Remplir les orbitales les plus basses, deux électrons chacune~;
        former la nouvelle matrice $\mathbf P$ et l'énergie.
  \item Reprendre à l'étape 3 jusqu'à ce que l'énergie et $\mathbf P$
        varient de moins qu'un seuil.
\end{enumerate}
\end{method}

\begin{omfigure}
\begin{tikzpicture}[font=\small, node distance=4mm,
  box/.style={draw=black!60, rounded corners=2pt, fill=omDef!6, align=center,
  text width=4.4cm, inner sep=4pt}]
  \node[box] (a) {intégrales $S$, $h$, $(\mu\nu|\lambda\sigma)$};
  \node[box, below=of a] (b) {estimation de $\mathbf P$};
  \node[box, below=of b] (c) {construire $\mathbf F(\mathbf P)$};
  \node[box, below=of c] (d) {résoudre $\mathbf F\mathbf C = \mathbf S\mathbf C\boldsymbol\varepsilon$};
  \node[box, below=of d] (e) {nouvelle $\mathbf P$, énergie $E$};
  \node[diamond, aspect=2.4, draw=black!60, fill=omThm!6, below=of e, inner sep=1pt] (f) {convergé~?};
  \node[box, below=of f, fill=omProp!8] (g) {orbitales, $\varepsilon$, $E$};
  \draw[->] (a) -- (b); \draw[->] (b) -- (c); \draw[->] (c) -- (d);
  \draw[->] (d) -- (e); \draw[->] (e) -- (f);
  \draw[->] (f) -- node[left] {oui} (g);
  \draw[->] (f.east) -- ++(1.6,0) |- node[pos=0.25, right] {non} (c.east);
\end{tikzpicture}

{\small La boucle du \omterm{def:b3:computational-chemistry:hartree-fock}{champ autocohérent} d'un calcul Hartree--Fock.}
\end{omfigure}

\begin{example}[\ce{H2} en base minimale]\label{ex:b3:computational-chemistry:h2}
% ledger: re:H2, janaf:H.dfH0, diat:H2.we, diat:H2.wexe
Avec une fonction STO-3G par atome, l'orbitale liante est fixée par la
symétrie,
$\sigma_g \propto \chi_A + \chi_B$, et le SCF ne demande qu'une étape. À $R =
\qty{1.4}{\bohr}$, le recouvrement vaut $S_{AB} = 0.6593$ et l'énergie
\qty{-1.1167}{\hartree}~; le minimum se situe à \qty{1.346}{\bohr}
(\qty{71.2}{pm},
contre \qty{74.1}{pm} mesurés) et \qty{-1.1175}{\hartree}. La figure
ci-dessous compare cette courbe à la vraie, construite à partir des
données spectroscopiques mesurées et de l'énergie de dissociation (une
courbe de Morse, \cref{ch:b3:rovibrational-spectroscopy}).
\end{example}

\begin{omfigure}
\begin{tikzpicture}
\begin{axis}[width=0.7\linewidth, height=6cm, xmin=0.8, xmax=10, ymin=-1.2,
  ymax=-0.45, axis lines=left, xlabel={$R$ / $a_0$}, ylabel={$E$ / $E_h$},
  legend cell align=left, legend style={font=\small, at={(0.04,0.98)}, anchor=north west, draw=none},
  tick label style={font=\small}, label style={font=\small}]
\addplot[omDef, thick] table[x=R, y=rhf] {figdata/out/bachelor-3-minimal-basis-hf-h2.dat};
\addlegendentry{RHF/STO-3G}
\addplot[black, thick, dashed] table[x=R, y=morse] {figdata/out/bachelor-3-minimal-basis-hf-h2.dat};
\addlegendentry{mesurée (forme de Morse)}
\addplot[gray, dotted, domain=0.8:10] {-1.0};
\addplot[omThm, dotted, domain=0.8:10] {-0.93316};
\node[font=\scriptsize, anchor=north] at (axis cs:8,-1.005) {deux atomes H exacts};
\node[font=\scriptsize, anchor=south, omThm] at (axis cs:8,-0.930) {deux atomes H en STO-3G};
\end{axis}
\end{tikzpicture}

{\small L'énergie de \ce{H2} en fonction de la distance internucléaire~:
Hartree--Fock restreint dans la base STO-3G (trait plein) et courbe
réelle (en tirets, forme de Morse construite à partir des $D_0$,
$\omega_e$, $\omega_ex_e$, $r_e$ mesurés). Près du
minimum, RHF est bon~; à grand $R$, il s'élève bien au-dessus de deux
atomes d'hydrogène.}
\end{omfigure}

\begin{proposition}[Pourquoi Hartree--Fock restreint échoue à la dissociation]\label{prop:b3:computational-chemistry:rhf-dissociation}
En \omterm{def:b3:computational-chemistry:basis}{base minimale}, la fonction d'onde RHF de \ce{H2} est, à la
normalisation près, $\sigma_g(1)\sigma_g(2) \propto \chi_A(1)\chi_B(2) + \chi_B(1)\chi_A(2)
+ \chi_A(1)\chi_A(2) + \chi_B(1)\chi_B(2)$~: à toute distance, à moitié
covalente et à moitié ionique (\ce{H+ H-}). À grand $R$, son énergie
tend vers la moyenne de deux atomes neutres et d'une paire d'ions, et
non vers celle de deux atomes.
\end{proposition}

\begin{proof}
Développons $(\chi_A + \chi_B)(1)(\chi_A + \chi_B)(2)$~: les deux termes
croisés placent un électron sur chaque atome, les deux termes carrés les
deux électrons sur le même atome, avec des poids égaux quel que soit
$R$. Le déterminant ne peut pas modifier ces poids.
\end{proof}

\begin{definition}[Corrélation électronique, énergie de corrélation]\label{def:b3:computational-chemistry:correlation}
La \emph{corrélation électronique}\index{correlation electronique@corrélation électronique} est la part de
l'évitement mutuel des électrons que ne décrit pas un déterminant unique
(dont les électrons de spins opposés se déplacent indépendamment dans le
champ moyen les uns des autres). L'\emph{énergie de
corrélation}\index{energie de correlation@énergie de corrélation} est
l'énergie exacte diminuée de l'énergie Hartree--Fock en base complète~;
elle est toujours négative.
\end{definition}

L'\omterm{def:b3:computational-chemistry:correlation}{énergie de corrélation} vaut environ \qty{-0.04}{\hartree} par paire
d'électrons, un pour cent d'une énergie totale, mais elle est comparable
à l'énergie d'une réaction~: les méthodes qui l'ajoutent (interaction de
configurations, théorie des perturbations, clusters couplés) coûtent bien
plus cher que Hartree--Fock, dont le coût croît déjà à peu près comme la
puissance quatrième du nombre de fonctions de base.

\begin{theorem}[Théorème de Koopmans]\label{thm:b3:computational-chemistry:koopmans}
Dans la théorie de Hartree--Fock, l'énergie nécessaire pour arracher un
électron à l'orbitale occupée $a$, toutes les autres orbitales restant
figées, vaut $-\varepsilon_a$. Les énergies d'ionisation mesurées sont
donc approximativement les opposées des énergies des orbitales~: c'est
le \emph{théorème de Koopmans}\index{theoreme de Koopmans@théorème de Koopmans}.
\end{theorem}

\begin{proof}
L'énergie d'un déterminant à couches complètes vaut $\sum_a2h_{aa} + \sum_{ab}(2J_{ab} -
K_{ab})$, et $\varepsilon_a = h_{aa} + \sum_b(2J_{ab} - K_{ab})$. Arrachons
un électron à $a$ sans modifier les orbitales~: les termes perdus sont
$h_{aa}$, les interactions de cet électron avec tous les autres,
$\sum_b(2J_{ab} - K_{ab}) -
J_{aa}$, plus $J_{aa}$ avec son ancien partenaire --- au total exactement
$\varepsilon_a$. Donc $E^+ - E = -\varepsilon_a$.
\end{proof}

\begin{example}[Ioniser \ce{H2}]\label{ex:b3:computational-chemistry:koopmans}
% ledger: ie4:H2
À \qty{1.4}{\bohr}, l'énergie de l'orbitale $\sigma_g$ en STO-3G vaut
\qty{-0.5782}{\hartree}~: Koopmans prévoit \qty{15.73}{eV}, contre
\qty{15.43}{eV} mesurés. Les orbitales figées surestiment l'énergie de
l'ion (qui se relaxerait), la corrélation manquante agit dans l'autre
sens~; pour les ionisations de valence, les deux erreurs se compensent
largement, et c'est pourquoi les spectres de photoélectrons du volume de
deuxième année se lisent à l'aide de diagrammes d'orbitales.
\end{example}

\section{Théorie de la fonctionnelle de la densité, et ce que dit un calcul}

La fonction d'onde de $N$ électrons dépend de $3N$ coordonnées~; la
\omterm{def:b3:computational-chemistry:dft}{densité électronique}, de trois.

\begin{definition}[Densité électronique, théorie de la fonctionnelle de la densité]\label{def:b3:computational-chemistry:dft}
La \emph{densité électronique}\index{densite electronique@densité électronique} $\rho(\mathbf r)$ est le
nombre d'électrons par unité de volume en $\mathbf r$, sommé sur tous les
électrons~:
$\int\rho\,\dd\tau = N$. La \emph{théorie de la fonctionnelle de la densité}\index{theorie de la fonctionnelle de la densite@théorie de la fonctionnelle de la densité}
(DFT) calcule l'énergie de l'état fondamental comme une fonctionnelle de
$\rho$. En pratique, $\rho = \sum_a|\phi_a|^2$ est construite à partir
d'\emph{orbitales de Kohn--Sham}\index{orbitales de Kohn--Sham}, les
orbitales d'électrons fictifs indépendants qui ont la même densité que
les électrons réels, et tous les effets d'échange et de corrélation sont
regroupés dans une \emph{fonctionnelle d'échange--corrélation}\index{fonctionnelle d'echange--correlation@fonctionnelle d'échange--corrélation} $E_{\mathrm{xc}}[\rho]$.
\end{definition}

\begin{theorem}[Hohenberg--Kohn]\label{thm:b3:computational-chemistry:hohenberg-kohn}
La densité de l'état fondamental d'un système d'électrons détermine le
potentiel extérieur (les noyaux) à une constante près, donc le
hamiltonien et toutes les propriétés de l'état fondamental~; et la
fonctionnelle énergie de la densité est minimale pour la vraie densité.
\end{theorem}
\admitted

La démonstration relève de cours plus avancés. Elle garantit qu'une
fonctionnelle exacte existe, non ce qu'elle est~: toute fonctionnelle
pratique est une approximation, choisie dans une hiérarchie --- densité
locale, corrections de gradient, fonctionnelles hybrides incorporant une
part d'échange Hartree--Fock --- et calibrée sur des données de
référence. Les équations de Kohn--Sham se résolvent par la même boucle
SCF que Hartree--Fock, pour un coût voisin, mais incluent la
corrélation~; c'est pourquoi la DFT est devenue l'outil de travail
quotidien de la chimie.

\begin{method}[Lire un calcul]\label{met:b3:computational-chemistry:read-calculation}
\begin{enumerate}
  \item Vérifier que le SCF a convergé et que l'\omterm{def:b3:computational-chemistry:pes}{optimisation de géométrie}
        s'est achevée avec des forces sous le seuil.
  \item Lancer un calcul de fréquences au même niveau~: des fréquences
        toutes réelles signalent un minimum~; une fréquence imaginaire,
        une structure de transition, dont le mode normal est le
        mouvement de franchissement de la barrière.
  \item Comparer ce qui est comparable~: des différences d'énergie entre
        structures calculées à un même niveau, jamais des énergies
        absolues issues de niveaux différents.
  \item Comparer à l'expérience en connaissant les erreurs typiques du
        niveau~: les fréquences harmoniques sont en général multipliées
        par un facteur d'échelle de 0.9 à 0.97~;
        les longueurs de liaison sont justes à environ 1 pm avec une
        fonctionnelle hybride et une base polarisée~; les énergies de
        réaction à quelques \unit{kJ/mol} au mieux.
\end{enumerate}
\end{method}

\begin{method}[Choisir un niveau de théorie]\label{met:b3:computational-chemistry:choose-level}
Géométries et fréquences de molécules organiques~: une fonctionnelle
hybride avec une \omterm{def:b3:computational-chemistry:split-valence}{base à valence séparée} polarisée. Interactions
faibles~: ajouter une correction de dispersion et des \omterm{def:b3:computational-chemistry:split-valence}{fonctions
diffuses}. Barrières de réaction et énergétique précise~: méthodes
corrélées de fonction d'onde sur des géométries DFT, autant que le
budget le permet. Grands systèmes (protéines, surfaces)~: une petite
région traitée par la mécanique quantique au sein d'un champ de forces
classique. Toujours tester le niveau choisi sur une molécule apparentée
dont la réponse est connue.
\end{method}

\begin{inthelab}[Une expérience numérique]
Un calcul se conduit comme une expérience et se consigne dans le cahier
de laboratoire~: le programme et sa version, la méthode, la base, les
seuils de convergence, la structure de départ, et les fichiers de sortie
conservés comme données brutes. Une séquence typique~: construire la
molécule, optimiser sa géométrie à un niveau modeste, calculer les
fréquences pour confirmer un minimum, puis calculer l'énergie à un niveau
plus élevé sur cette géométrie.
\end{inthelab}

\begin{history}[Une molécule connue d'abord par le calcul]
% ledger: hist4:HeHp.2019
\ce{HeH+} fut produit au laboratoire en 1925. Ses propriétés furent
calculées en détail à partir des années 1930, et l'on prédit qu'il
devait être abondant dans certains gaz astrophysiques~; mais sa première
raie de rotation se situe dans l'infrarouge lointain, absorbé par
l'atmosphère. En 2019, un observatoire aéroporté le détecta dans la
nébuleuse planétaire NGC 7027, à la fréquence que les calculs et les
spectres de laboratoire avaient prévue.
\end{history}

\section{Exercices}

\begin{exercise}[$\star$]\label{exo:b3:computational-chemistry:1}
Pour l'atome d'hydrogène, prendre la fonction d'essai $\eu^{-\alpha r^2}$.
En unités atomiques, $E(\alpha) = \frac32\alpha - 2\sqrt{2\alpha/\pi}$.
Trouver le meilleur $\alpha$ et l'énergie correspondante~; comparer à
$-\frac12E_h$.
\end{exercise}

\begin{exercise}[$\star$]\label{exo:b3:computational-chemistry:2}
Combien de fonctions de base un calcul sur le benzène utilise-t-il dans
la base STO-3G~? Et en 6-31G(d), avec six fonctions $d$ cartésiennes par
carbone et deux fonctions $s$ par hydrogène~?
\end{exercise}

\begin{exercise}[$\star$]\label{exo:b3:computational-chemistry:3}
Expliquer pourquoi \ce{H2} et \ce{D2} ont la même longueur de liaison et
la même \omterm{def:b3:quantum-model-systems:oscillator}{constante de force}, mais des nombres d'onde de vibration
différents. Quel rapport attend-on~?
\end{exercise}

\begin{exercise}[$\star$]\label{exo:b3:computational-chemistry:4}
Un calcul de fréquences sur une structure de formule \ce{C2H5F} donne 17
fréquences réelles et une fréquence imaginaire, $487\iu$~\unit{cm^{-1}}.
De quel type de \omterm{def:b3:computational-chemistry:pes}{point stationnaire} s'agit-il~? Combien de fréquences
attendait-on au total~?
\end{exercise}

\begin{exercise}[$\star\star$]\label{exo:b3:computational-chemistry:5}
% ledger: ie:He, ie4:He2
À l'aide de $E(\zeta) = \zeta^2 - \frac{27}{8}\zeta$, calculer la
meilleure énergie variationnelle de l'hélium en \unit{\hartree} et en eV,
et son erreur par rapport à la valeur exacte
$\qty{-2.9034}{\hartree}$. Quelle fraction de l'énergie totale
l'erreur représente-t-elle~?
\end{exercise}

\begin{exercise}[$\star\star$]\label{exo:b3:computational-chemistry:6}
Deux fonctions de base donnent $H_{11} = \qty{-1.0}{\hartree}$, $H_{22} =
\qty{-0.5}{\hartree}$, $H_{12} = \qty{-0.2}{\hartree}$, $S_{12} = 0.3$.
Résoudre l'équation séculaire. La plus petite racine est-elle
au-dessous de $H_{11}$~?
\end{exercise}

\begin{exercise}[$\star\star$]\label{exo:b3:computational-chemistry:7}
Un calcul Hartree--Fock dure 10 minutes pour le benzène en STO-3G. En
supposant un coût proportionnel à la puissance quatrième du nombre de
fonctions de base, estimer la durée en 6-31G(d).
\end{exercise}

\begin{exercise}[$\star\star$]\label{exo:b3:computational-chemistry:8}
% ledger: ie4:H2
L'énergie de l'orbitale STO-3G de \ce{H2} à \qty{1.4}{\bohr} vaut
\qty{-0.5782}{\hartree}. Donner l'estimation de Koopmans de l'énergie
d'ionisation en eV et la comparer à \qty{15.43}{eV}. Citer deux raisons
de leur écart, avec leur sens.
\end{exercise}

\begin{exercise}[$\star\star$]\label{exo:b3:computational-chemistry:9}
En STO-3G, l'énergie d'un atome d'hydrogène vaut \qty{-0.4666}{\hartree}
et l'énergie RHF de \ce{H2} à \qty{10}{\bohr} vaut \qty{-0.5960}{\hartree}.
De combien l'énergie RHF dépasse-t-elle celle de deux atomes à cette
distance, en eV~? Expliquer à l'aide des termes ioniques.
\end{exercise}

\begin{exercise}[$\star\star\star$]\label{exo:b3:computational-chemistry:10}
Pour $\phi = c_1\chi_1 + c_2\chi_2$ avec des fonctions réelles, écrire
$E(c_1,c_2)$ et établir les deux équations séculaires en posant
$\partial E/\partial c_1 =
\partial E/\partial c_2 = 0$.
\end{exercise}

\begin{exercise}[$\star\star\star$]\label{exo:b3:computational-chemistry:11}
% ledger: aw:H, diat:H2.we
L'énergie RHF/STO-3G de \ce{H2} vaut $-1.116871$, $-1.117501$ et
$\qty{-1.116714}{\hartree}$ à $R = 1.30$, 1.35 et \qty{1.40}{\bohr}.
Estimer la \omterm{def:b3:quantum-model-systems:oscillator}{constante de force} par différences finies, la convertir en
\unit{N/m}
($\qty{1}{\hartree\per\bohr\squared} = \qty{1556.9}{N/m}$), et calculer le
nombre d'onde harmonique. Comparer à $\tilde\omega_e = \qty{4401}{cm^{-1}}$.
\end{exercise}

\begin{exercise}[$\star\star\star$]\label{exo:b3:computational-chemistry:12}
On doit calculer la différence d'énergie entre deux conformères d'un
sucre, dont on sait qu'ils diffèrent de quelques \unit{kJ/mol}, l'un
d'eux possédant une liaison hydrogène intramoléculaire. Choisir et
justifier une méthode, une base et les vérifications à effectuer.
\end{exercise}

\section{Problème~: HeH\texorpdfstring{$^+$}{+}, la première molécule}

\begin{problem}[{Problème du week-end --- un calcul Hartree--Fock en base
minimale de l'ion hydrure d'hélium~: la base, l'estimation de cœur,
l'orbitale convergée, et la liaison d'un proton à l'hélium}]
\label{pb:b3:computational-chemistry:1}
% ledger: bse:He.sto3g.a1, bse:He.sto3g.a2, bse:He.sto3g.a3, bse:H.sto3g.a1, pa:He, ie:He, ie4:He2
On calcule \ce{HeH+} à $R = \qty{1.4632}{\bohr}$ dans la base STO-3G
(exposants de He
6.3624, 1.1589, 0.31365~; exposants de H 3.4253, 0.62391, 0.16886~; les
trois mêmes coefficients). La fonction 1 est centrée sur He, la fonction
2 sur H. Le programme donne (en unités atomiques)
\[
\mathbf S = \begin{pmatrix}1 & 0.5368\\ 0.5368 & 1\end{pmatrix},\quad
\mathbf h = \begin{pmatrix}-2.5983 & -1.4318\\ -1.4318 & -1.7318\end{pmatrix},\quad
\mathbf F_{\mathrm{final}} = \begin{pmatrix}-1.5902 & -1.0610\\ -1.0610 & -0.8340
\end{pmatrix},
\]
l'orbitale occupée $\phi = 0.8766\chi_1 + 0.2025\chi_2$, les énergies des
orbitales
$-1.6328$ et \qty{-0.1725}{\hartree}, et les énergies SCF successives
$-2.7978$, $-2.8404$,
$-2.8418$, $-2.8418$. Un atome d'hélium dans la même base a pour énergie
\qty{-2.8078}{\hartree}. $\qty{1}{\hartree} = \qty{27.211}{eV} =
\qty{2625.5}{kJ/mol}$.

\medskip\noindent\textbf{Partie I --- La base.}
\begin{enumerate}
  \item Que signifie «~STO-3G~»~?
  \item Pourquoi les exposants de l'hélium sont-ils plus grands que ceux
        de l'hydrogène~? Vérifier que leur rapport vaut $(1.69/1.24)^2$.
  \item Combien y a-t-il de fonctions de base, d'électrons et d'orbitales
        occupées~?
  \item Calculer l'énergie de répulsion nucléaire.
  \item Interpréter $S_{12} = 0.5368$.
  \item Pourquoi $h_{11}$ est-il plus bas que $h_{22}$~?
\end{enumerate}

\medskip\noindent\textbf{Partie II --- L'estimation de cœur.}
\begin{enumerate}[resume]
  \item Écrire l'équation séculaire $\det(\mathbf h - \varepsilon\mathbf S) = 0$
        sous forme d'une équation du second degré en $\varepsilon$.
  \item La résoudre.
  \item Pourquoi n'est-ce qu'un point de départ~?
  \item Donner l'expression de la matrice de Fock en fonction de $\mathbf h$,
        de la matrice densité et des intégrales biélectroniques.
  \item Quelles interactions entre électrons $\mathbf F$ contient-elle,
        qui manquent à $\mathbf h$~?
  \item Expliquer pourquoi la procédure doit être itérée.
\end{enumerate}

\medskip\noindent\textbf{Partie III --- L'orbitale convergée.}
\begin{enumerate}[resume]
  \item Vérifier que $\phi$ est normée.
  \item Les populations sont les éléments diagonaux de $P\mathbf S$, $2c_1^2 + 2c_1c_2S$
        sur He et $2c_2^2 + 2c_1c_2S$ sur H. Les calculer.
  \item Où se trouve la charge positive~? La molécule est-elle mieux
        décrite par \ce{He + H+} ou par \ce{He+ + H}~?
  \item Estimer, à l'aide du \omterm{thm:b3:computational-chemistry:koopmans}{théorème de Koopmans}, l'énergie nécessaire
        pour arracher un électron à \ce{HeH+}.
  \item Calculer l'énergie électronique $E_{\mathrm{el}} = E - V_{\mathrm{nn}}$.
  \item Combien de cycles a-t-il fallu pour converger à
        \qty{e-4}{\hartree} près~? Pourquoi l'énergie a-t-elle diminué à
        chaque cycle~?
  \item Avec des exposants de l'hélium mis à l'échelle de $\zeta = 2.0925$
        au lieu de 1.69, le même programme donne \qty{-2.8607}{\hartree}.
        Quelle base est la meilleure, et pourquoi peut-on l'affirmer~?
  \item Que représente l'énergie de l'orbitale virtuelle, \qty{-0.1725}{\hartree}~?
\end{enumerate}

\medskip\noindent\textbf{Partie IV --- La liaison d'un proton.}
\begin{enumerate}[resume]
  \item Quelle est l'énergie d'un proton nu~? Celle de \ce{He + H+} très
        éloignés, dans cette base~?
  \item Calculer l'énergie libérée par \ce{He + H+ -> HeH+} à ce niveau,
        en \unit{kJ/mol}.
  \item L'affinité protonique mesurée de l'hélium vaut
        \qty{177.8}{kJ/mol}. Commenter.
  \item L'énergie exacte de l'hélium est l'opposé de la somme de ses
        énergies d'ionisation, 24.5874 et \qty{54.4178}{eV}. La calculer
        en \unit{\hartree}, ainsi que l'erreur de l'atome d'hélium en
        STO-3G.
  \item Pourquoi une erreur de cette taille peut-elle encore donner des
        géométries moléculaires utiles~?
  \item Énoncer le résultat~: l'énergie totale RHF/STO-3G de \ce{HeH+} à
        \qty{1.4632}{\bohr}.
\end{enumerate}
\end{problem}
